% page 299 of the facsimile (image p-298 of the scan)
% block 167.6 x 242.0 mm ; left margin 34.4 mm
\pgc{12.700mm}{0.000mm}{
\l{komplikar. (\sou{trans, per}).\cel{2}Inkombrar per elementi multa-sorta.}
\l{\cel{4}- DEFIS.}
\l{}
\l{\sou{komplimentar}. (\sou{trans., pri, pro}). I. Diskursar homaje di (ulu),}
\l{\cel{4}respekt-exprese di (ulu), en okaziono solena. - II. Dicar}
\l{\cel{4}(ad ulu) vorti qui laudas lu. - III. Dicar politeso-vorti.}
\l{\cel{4}- DEFIRS.}
\l{}
\l{\sou{komplotar}. (\sou{netrans.})\cel{2}I. Projetar, kombine da plura personi,}
\l{\cel{4}kontre la vivo, kontre la sekureso, di ulu. - II.\cel{1}\sou{(Fami}-}
\l{\cel{4}\sou{liare}). Kombinar ulo kontre ulu. - DEFRS.}
\l{}
\l{\sou{kompostar}. (\sou{trans.)}\cel{3}Asemblar la tipi imprimala por reproduk-}
\l{\cel{4}tar la texto imprimenda. - EFIS.}
\l{}
\l{\sou{kompoto}. Entremeso sukroza, ek frukti integra o quarima, koqui-}
\l{\cel{4}ta kun sukro, aquo o vino. - DEFS.}
\l{}
\l{\sou{kompozar}. (\sou{trans.})\cel{2}I. Formacar (ensemblo) per la asemblo o la}
\l{\cel{4}kombino di elementi diversa. - II. Facar verko skriptala,}
\l{\cel{4}per kombinar olua elementi, olua parti. - DEFIS.}
\l{}
\l{\sou{komprar}. (\sou{trans., po, por, de)}. - Aquirar po pekunio-quanto. IS.}
\l{}
\l{\sou{komprenar}. (\sou{trans.})\cel{2}Embracar per la penso (la senco di ulo,}
\l{\cel{4}la naturo di ulo, la motivo di ulo). - EFIS.}
\l{}
\l{\sou{kompresar}. (\sou{trans.)}\cel{2}Reduktar a volumino min granda, per presar}
\l{\cel{4}o klemar. - DEFIRS.}
\l{}
\l{\sou{kompromisar}. (\sou{trans.})\cel{2}Igar en situeso danjeroza (pri lua}
\l{\cel{4}interesti, reputeso, digneso). - DEFIRS.}
\l{}
\l{\sou{kompunciono}.\cel{2}I. (\sou{teol. kato}l.)Tristeso pia, qua devenas de}
\l{\cel{4}la koncio di nia nedigneso e de la sufro ofensir Deo. -}
\l{\cel{4}II. Mieno grava, rekolecala. - EFIS.}
\l{}
\l{kompunda. (\sou{tekn.)}\cel{3}Dicesas pri ula organi od aparati qui esas}
\l{\cel{4}\textquotedbl{}asociita\textquotedbl{},kombinita reciproke.\cel{2}- EF.}
\l{}
\l{\sou{komto.}\cel{2}I. En la epoki finala di la imperio Romana, e komence di}
\l{\cel{4}la mez-epoko : oficiro di la suvereno-palaco, mayoro, gu-}
\l{\cel{4}vernisto di teritorio-regiono. - II. Lor la feudismo : sinio-}
\l{\cel{4}ro di feudo, qua esis, en la hierarkio, nemediate sub la}
\l{\cel{4}dukio. - III. En la hierarkio ek la nobelo-tituli : la persono}
\l{\cel{4}qua esas nemediate dop la markezo. - EFIS.}
\l{}
\l{komuna.\cel{3}Qua apartenas kune ad plura. - DEFIRS.}
\l{}
\l{komuniar. (\sou{netrans.}) (\sou{religio kristana}).\cel{2}Recevar la sakramento}
\l{\cel{4}\textquotedbl{}eukaristio\textquotedbl{}. - DEFIS.}
\l{}
\l{komunikar. (trans., ad,\cel{1}\sou{kun}). I. Igar (kozo quan onu posedas)}
\l{\cel{4}komuna ad altru per transmisar olu a lu. II. Relateskar}
\l{\cel{4}ideale, interestale e c. kun ulo. III. Inter-relatar per}
\l{\cel{4}paseyo. - DEFIS.}
}
